\documentclass[11pt]{article}
\usepackage[margin=1in]{geometry}
\usepackage{booktabs}
\usepackage{url}
\title{Why two 8-m/s encounters stop at 8584c47}
\date{Prepared October 4, 2026}
\begin{document}
\maketitle

\textbf{Scope.} The rerun in \path{CONTROLLER_CAMPAIGN_20261004.tex} stopped
in two cases with \texttt{noOptimizationSolution
(linearizationAccuracyNotReached)}: 8-m/s turningCrossing at hold 20 and
8-m/s curvedCrossing at hold 219. Both cases recovered at 91a35f6. This
report identifies the change behind each stop and the mechanism. No
production source was changed. All interventions were made in temporary
worktrees or an instrumented scratch copy; the diff is
\path{CONTROLLER_FAILURE_DIAGNOSIS_20261004/instrumentation.diff}.

\section{What the stop means}
Each hold, \texttt{solvePredictiveControl} linearizes the bicycle model about
an anchor trajectory and solves a two-stage SOCP. The first stage minimizes
PCBF slack and the second minimizes CLF slack. The first hold of the result
is then replayed through the nonlinear model. The agreement ratio is the
largest of three scaled errors: pose error over 1 cm, state error over
$[5,3,1.5]\cdot 1\%$, and CLF prediction error over 0.01 (scaled by
$\max(1,W_0)$). A ratio above 1 rejects the plan. Two recoveries are then
tried. First, a damped step is taken, but only if the zero-correction anchor
is feasible. Second, one fresh potential-field seed is solved. There are at
most two linearizations per hold. If neither recovery agrees, the hold
issues no command and the experiment stops.

\section{Isolation}
Both failing frames replay deterministically from the saved snapshots.
Bisection over the controller-affecting commits
(\path{isolation-results.txt}) gives the following:

\begin{center}
\begin{tabular}{lcc}
\toprule
Source & turningCrossing & curvedCrossing \\
\midrule
91a35f6, 838848b, 388dcfc & recovered & recovered \\
aa85322, 73abcad & recovered & stops (hold 185) \\
8584c47 (HEAD) & stops (hold 20) & stops (hold 219) \\
\midrule
HEAD, agreement over encounter restored & recovered & stops (hold 183) \\
HEAD, dual instead of support linearization & stops (hold 20) & stops (hold 210) \\
HEAD, encounter range 30 m & stops (hold 20) & recovered (hold 377) \\
aa85322, encounter range 30 m & -- & recovered (hold 196) \\
\bottomrule
\end{tabular}
\end{center}

The two stops have separate causes. turningCrossing follows the 8584c47
reduction of the agreement check from every in-range node to the first
hold. curvedCrossing follows the aa85322 increase of the encounter range from
30 m to a 50 m reference-position radius, which also permits re-entry.

\section{turningCrossing: first-hold agreement lets the plan drift}
After 8584c47 only the issued hold must agree. In the HEAD trace, the
first-hold pose error stays at about 2~mm, so the input trust scale grows
every hold, from 0.125 to 0.751 by 0.55 s. The unchecked remainder of the
plan grows with it: its pose error relative to the nonlinear model reaches
5.4~cm, five times the old 1-cm tolerance. Each next hold restarts from that
shifted plan. Measured against the new linearization, the shifted plan
violates the constraint set by 0.11. At 1.0~s, 0.26~s before the baseline
collision, that violation jumps to 0.616. The optimizer then needs a full
trust-radius steering change (0.224 to 0.149 rad). The resulting yaw-rate
prediction is off by 0.034 rad/s, so the state ratio is 2.24. Damping is
impossible because the anchor itself is infeasible. The fresh
potential-field seed returns no numerical PCBF result at trust scales 0.125
and 2. The hold therefore stops.

With the encounter-wide check restored (and everything else at HEAD), the
trust scale stays at or below 0.29. The full-plan error stays at or below
0.8~cm, the shifted-plan violation stays at or below 0.015, and the case
recovers.

\section{curvedCrossing: the target returns during recovery}
The target is an NRMM motion with constant sideslip 0.04~rad and rear-axle
distance 1.6~m. Its turning radius is therefore
$1.6/\sin 0.04 = 40.0$~m, regardless of its acceleration. The fixture keeps
this circling for all time, so after the first crossing the target comes
back. With range 30 m at HEAD, the ego meets it again at a 5.5~m center
distance at 12.55~s. The CLF there rises from 0.8 to 371, a second avoidance,
and the case recovers at 18.85~s.

With the 50 m radius, the returning target stays in the constraint set while
the ego is still finishing its first recovery. At the failing hold
($t=10.95$~s) the target is 25~m away. The horizon of about 4.1~s, however,
predicts rectangle overlap along the anchor 1.6 to 1.8~s ahead (row value
$-0.98$~m). From 10.45~s the CLF stops decreasing, and the CLF slack grows
from 0.01 to 1.01. The plans now need large corrections, and the quadratic
CLF prediction loses accuracy: predicted 2.537 against nonlinear 2.281
(ratio 15.0), and 7.4 after the fresh seed. The best damped trial reaches
1.155. Replaying the failing hold without the per-node collision rows lets
the hold succeed. Removing only the terminal target rows does not.

At 91a35f6 this case was recorded as recovered at 9.75~s, before the target
returns. The experiment stops at the first confirmed recovery, so the second
encounter was never exercised in that record.

\section{Interpretation and options}
Neither stop is a sampled collision. Both are failures to return a
nonlinearly consistent command. They share the same final step: two
linearizations per hold and no fallback when both disagree.
\begin{itemize}
\item turningCrossing is a regression from the 8584c47 relaxation, which was
made for the noisy-estimator campaign. Restoring the encounter-wide check
fixes this case, but its effect on the noisy campaign has not been tested.
Limiting trust-scale growth by the unchecked plan error is an untested
alternative.
\item curvedCrossing is a real second encounter that the 50-m range exposes
earlier, combined with a CLF agreement test that is strict while CLF slack is
positive. A shorter range hides the return until the ego has settled. The
campaign's stop-at-first-recovery criterion also hid it before.
\end{itemize}
None of these options was implemented or validated here.

\section{Reproduction}
Inputs are the saved failure snapshots and traces in
\path{~/.cache/collisionAvoidance/controller-campaign-20261004/}. The
scripts in \path{CONTROLLER_FAILURE_DIAGNOSIS_20261004/} are stored as
\path{.m.txt}: \texttt{bisectCase}, \texttt{bisectCase2} and
\texttt{traceCase} run cases in a worktree; \texttt{replayFailure} with
\texttt{runDiag1}--\texttt{6} replays a failing hold in the instrumented
copy. Every run uses exact target observations, the \texttt{ode45} replay,
a 1500-hold cap, a 1-s recovery dwell and the 0.20~m margin.
\end{document}
